% An \index after other forms of minted's listings, each in a section of
% its own, with makeidx and without hyperref: \inputminted, \mint and
% \mintinline, an environment from \newminted, a command from
% \newmintedfile, a listing then a footnote, tcolorbox's tcblisting and
% \tcbinputlisting with minted, and minted's listing float. Each number
% links to the place of its \index, right after its listing.
\documentclass{article}
\usepackage{minted,makeidx,tcolorbox}
\tcbuselibrary{minted}
\newminted{python}{linenos}
\newmintedfile[pyfile]{python}{linenos}
\makeindex
\begin{document}
\section{Input}
\inputminted[linenos]{python}{test-idxminted-code.py}
After the input listing\index{one}.
\section{Inline}
\mint{python}|total = total + 1|
Then \mintinline{python}{total = 2} inline, and after both\index{two}.
\section{Environment}
\begin{pythoncode}
total = total + 0
total = total + 1
total = total + 2
total = total + 3
total = total + 4
total = total + 5
total = total + 6
total = total + 7
total = total + 8
total = total + 9
total = total + 10
total = total + 11
total = total + 12
total = total + 13
total = total + 14
total = total + 15
total = total + 16
total = total + 17
total = total + 18
total = total + 19
\end{pythoncode}
After the pythoncode listing\index{three}.
\section{File}
\pyfile{test-idxminted-code.py}
After the pyfile listing\index{four}.
\section{Footnote}
\begin{minted}[linenos]{python}
total = total + 0
total = total + 1
total = total + 2
total = total + 3
total = total + 4
total = total + 5
total = total + 6
total = total + 7
total = total + 8
total = total + 9
total = total + 10
total = total + 11
total = total + 12
total = total + 13
total = total + 14
total = total + 15
total = total + 16
total = total + 17
total = total + 18
total = total + 19
\end{minted}
Text with a note.\footnote{In the footnote\index{five}.}
\section{Box}
\begin{tcblisting}{listing engine=minted, minted language=python, listing only}
total = total + 0
total = total + 1
total = total + 2
total = total + 3
total = total + 4
total = total + 5
total = total + 6
total = total + 7
total = total + 8
total = total + 9
total = total + 10
total = total + 11
total = total + 12
total = total + 13
total = total + 14
total = total + 15
total = total + 16
total = total + 17
total = total + 18
total = total + 19
\end{tcblisting}
After the box\index{six}.
\section{Input box}
\tcbinputlisting{listing engine=minted, minted language=python, listing only,
  listing file=test-idxminted-code.py}
After the input box\index{seven}.
\section{Float}
\begin{listing}[h]
\begin{minted}[linenos]{python}
total = total + 0
total = total + 1
total = total + 2
total = total + 3
total = total + 4
total = total + 5
total = total + 6
total = total + 7
total = total + 8
total = total + 9
total = total + 10
total = total + 11
total = total + 12
total = total + 13
total = total + 14
total = total + 15
total = total + 16
total = total + 17
total = total + 18
total = total + 19
\end{minted}
\caption{A listing in a float}
\end{listing}
After the float\index{eight}.
\printindex
\end{document}
